\section{Discusión}

\textbf{Respuesta a la pregunta.} Un modelo con rezagos, medias móviles y calendario reduce a la mitad el error relativo de la predicción a 24 horas frente a la referencia estacional semanal: en 2019, el Random Forest cíclico obtiene un NMAE macro de 0,0249 frente a 0,0505 (50,7\,\% menos; IC 95\,\% de 45,4\,\% a 54,9\,\%), con mejoras de 47\,\% a 54\,\% en cada país, y en el periodo de desarrollo (2015--2018) la mejora fue similar (49,6\,\%) y consistente en las tres ventanas. Buena parte de la señal ya está en el patrón semanal (\texttt{lag\_168h} concentra 65\,\% de la importancia); los modelos agregan sobre todo su corrección con el nivel reciente.

\textbf{¿Justifica la ganancia la complejidad?} En 2019, Ridge mejora a la referencia 18,8\,\%, el árbol podado 34,1\,\% y el Random Forest 50,7\,\%. La ventaja del bosque sobre el árbol (25,1\,\%; IC de 22,6\,\% a 27,7\,\%) y sobre Ridge (39,2\,\%; IC de 36,4\,\% a 42,1\,\%) es clara en los intervalos por bloques, aunque condicionales al ajuste y a un solo año. Su costo es menor interpretabilidad, 200 árboles con millones de nodos, ajustes de segundos (8,5 s el modelo final) y mayor dependencia de la historia (con solo el último año de entrenamiento el error sube 8,3\,\%). Con reajuste periódico ese costo es razonable; si se requiriera un modelo simple y auditable, el árbol podado conserva dos tercios de la mejora (34,1\,\% frente a 50,7\,\%). El calendario cíclico agrega una ganancia pequeña pero distinguible del ruido (4,8\,\%; IC de 3,1\,\% a 6,7\,\%).

\textbf{Límites de generalización.}
\begin{itemize}
    \item Los resultados valen para Alemania, España y Francia en 2015--2019, no para otros países ni para cambios de régimen; los árboles no extrapolan niveles fuera del rango de entrenamiento. La evaluación final usa un solo año, y los intervalos no incluyen la variabilidad entre años ni la de reajustar.
    \item Se predice de forma sucesiva a horizonte fijo, no emitiendo todo el día siguiente a una hora de mercado. No hay temperatura ni feriados: los peores días de 2019 son feriados (hasta 11\,200 MW de MAE diario), y en alta demanda el modelo subestima (sesgo de $-456$ a $-920$ MW).
    \item La importancia de variables es asociación predictiva, no causalidad. Las cifras del periodo de desarrollo (2015--2018) incluyen el efecto de elegir hiperparámetros con las mismas ventanas; la estimación imparcial es la de 2019. La versión cíclica usa hiperparámetros buscados con indicadores, y la grilla de poda y la de \texttt{max\_features} del bosque tienen sus mejores valores en el borde (Métodos).
    \item Las ausentes se tratan con casos completos y no son homogéneas entre países (Francia pierde 2,5\,\% de las filas); evaluar en fechas comunes cambia poco el NMAE, pero no se descarta que sean sistemáticas.
\end{itemize}

\textbf{Condiciones para usar el modelo.} Requiere las últimas 168 horas de demanda sin faltantes y el calendario local del instante objetivo. Debe reentrenarse con datos recientes y toda la historia, y vigilarse frente a la referencia semanal. No se recomienda confiar en él para feriados ni para el día siguiente a uno prolongado (por ejemplo, el 2 de enero), donde el error diario puede ser varias veces el habitual, ni para dimensionar capacidad en demanda extrema, que tiende a subestimar. Un calendario de feriados sería la primera mejora a evaluar.
